% ==============================================================================
% ADVANCED CONTROL LAW: NEW RNN, 3D MPC AND 3D EKF
% ==============================================================================

% ------------------------------------------------------------------------------
% NEW RNN MODEL
% ------------------------------------------------------------------------------
\begin{frame}{New RNN model}
    \framesubtitle{Trained on flight controller + drone simulations}

    \begin{columns}[T,onlytextwidth]
        \begin{column}{0.5\textwidth}
            \begin{block}{What the RNN learns}
                \centering
                \resizebox{\linewidth}{!}{%
                \begin{tikzpicture}[
                    >=Stealth, thick,
                    block/.style={draw, rectangle, minimum height=0.8cm,
                        minimum width=1.6cm, align=center, fill=blue!5,
                        rounded corners=2pt, font=\scriptsize}
                ]
                    \node[block] (fc) {\textbf{Flight}\\\textbf{controller}};
                    \node[block, right=1.0cm of fc] (drone) {\textbf{Drone}\\\tiny (Simulink)};
                    \draw[->] (fc) -- node[above, font=\tiny] {$U_{1..4}$} (drone);
                    \draw[->] ([xshift=-0.9cm]fc.west)
                        node[left, font=\scriptsize, align=center]
                        {$z_{\mathrm{ref}}, \phi_{\mathrm{ref}},$\\
                         $\theta_{\mathrm{ref}}, \psi_{\mathrm{ref}}$}
                        -- (fc.west);
                    \draw[->] (drone.east) -- ++(0.9,0)
                        node[right, font=\scriptsize] {$a_x, a_y, z$};
                    \draw[->, red!70!black] (drone.south) -- ++(0,-0.4) -|
                        node[pos=0.25, below, font=\tiny] {attitude, altitude} (fc.south);
                    \node[draw, dashed, gray, rounded corners, inner sep=6pt,
                        fit={(fc) (drone) ([yshift=-0.55cm]drone.south)
                             ([yshift=0.3cm]fc.north)}] (box) {};
                    \node[gray, font=\tiny\itshape, anchor=north]
                        at (box.north) {system identified by the RNN};
                \end{tikzpicture}%
                }
            \end{block}

            \begin{block}{Model and training}
                \footnotesize
                \begin{itemize}\setlength{\itemsep}{0pt}
                    \item Elman RNN, 32 hidden units, 1.3\,k parameters
                    \item 4 inputs $\rightarrow$ 3 outputs $(a_x, a_y, z)$
                    \item 3000 trajectories of 1941 steps, 70\,\% for training
                    \item Adam, $\mathrm{lr}=10^{-3}$, 100 epochs, batch 32
                    \item Test RMSE: $a_x$ 0.24, $a_y$ 0.23\,m/s$^2$,
                    $z$ 0.048\,m
                \end{itemize}
            \end{block}
        \end{column}

        \begin{column}{0.47\textwidth}
            \begin{alertblock}{Why a new model?}
                \footnotesize
                \begin{itemize}\setlength{\itemsep}{0pt}
                    \item The first RNN only described the $z$ axis
                    ($u \rightarrow a_z$).
                    \item 3D control needs $x$ and $y$, but raw motor data is
                    unusable (unstable drone).
                    \item With the flight controller inside, the identified
                    system is stable, with bounded physical inputs.
                    \item The MPC sends references to the flight
                    controller, like an autopilot.
                \end{itemize}
            \end{alertblock}

            \begin{block}{Choice of the outputs}
                \footnotesize
                \begin{itemize}\setlength{\itemsep}{0pt}
                    \item $z$ is held by the altitude loop: learned
                    directly, no double integration (lesson from $a_z$).
                    \item No position loop on $x, y$: learn $a_x, a_y$, then
                    integrate.
                \end{itemize}
            \end{block}
        \end{column}
    \end{columns}
\end{frame}

% ------------------------------------------------------------------------------
% 3D MPC
% ------------------------------------------------------------------------------
\begin{frame}{3D MPC}
    \framesubtitle{Position control through the flight controller}

    \begin{columns}[T,onlytextwidth]
        \begin{column}{0.52\textwidth}
            \begin{block}{Cost function minimized}
                \footnotesize
                \setlength{\abovedisplayskip}{0pt}\setlength{\belowdisplayskip}{2pt}
                \[
                    J(U) = \sum_{i=1}^{N} Q\,\big\|\hat{p}_{k+i} - r_{k+i}\big\|^2
                    + R\,\big\|u_{k+i-1} - u_{k+i-2}\big\|^2
                \]
                {\scriptsize\raggedright
                $\hat{p}=(\hat{x}, \hat{y}, \hat{z})$: position predicted
                with the RNN\\
                $u = (z_{\mathrm{ref}}, \phi_{\mathrm{ref}},
                \theta_{\mathrm{ref}}, \psi_{\mathrm{ref}})$: references
                sent to the flight controller\\
                $r$: 3D reference trajectory\par}
            \end{block}

            \begin{block}{Prediction over the horizon}
                \footnotesize
                \setlength{\abovedisplayskip}{0pt}\setlength{\belowdisplayskip}{2pt}
                \begin{align*}
                    (a_{x}, a_{y}, z)_{k} &= \mathrm{RNN}(h_{k-1}, u_{k-1})\\
                    v_{k} &= v_{k-1} + a_{k}\,\Delta t \quad (x, y)\\
                    p_{k} &= p_{k-1} + v_{k}\,\Delta t \quad (x, y)
                \end{align*}
                {\scriptsize $z$ is read directly from the RNN output.\par}
            \end{block}
        \end{column}

        \begin{column}{0.45\textwidth}
            \begin{block}{Constraints}
                \footnotesize
                $u_{\min} \le u \le u_{\max}$: the range seen in training
                \begin{itemize}\setlength{\itemsep}{0pt}
                    \item $z_{\mathrm{ref}} \in [0.5,\ 2.5]$\,m
                    \item $\phi_{\mathrm{ref}}, \theta_{\mathrm{ref}}
                    \in [-0.3,\ 0.3]$\,rad
                    \item $\psi_{\mathrm{ref}} \in [-0.5,\ 0.5]$\,rad
                \end{itemize}
                {\scriptsize Outside this range the RNN would extrapolate.\par}
            \end{block}

            \begin{exampleblock}{Settings}
                \footnotesize
                \begin{itemize}\setlength{\itemsep}{0pt}
                    \item 50 Adam steps per sample ($\mathrm{lr}=0.05$),
                    warm start
                    \item $T_e = 0.05$\,s, $Q = R = 1$
                    \item $N=20$ (1\,s): with $N=10$ (0.5\,s) the MPC diverges
                \end{itemize}
            \end{exampleblock}
        \end{column}
    \end{columns}
\end{frame}
